%%% historicity-journal.tex — Journal-length compressed version
%%% Compile: pdflatex historicity-journal && bibtex historicity-journal
%%%          && pdflatex historicity-journal (twice)
%%%
%%% Target: ~12,000 words, submittable to History of Religions,
%%% Journal for the Study of the Historical Jesus, The Journal of
%%% Religion, or similar venues. Full-length paper and two
%%% supplementary documents are cited as supporting material for
%%% evidence the journal article summarises rather than develops.

\documentclass[12pt,oneside]{article}
\usepackage[letterpaper,margin=1in,headheight=14pt]{geometry}
\usepackage[T1]{fontenc}
\usepackage[utf8]{inputenc}
\usepackage{amsmath,amssymb}
\usepackage{setspace}
\onehalfspacing
\usepackage{titlesec}
\usepackage{fancyhdr}
\pagestyle{fancy}
\fancyhf{}
\fancyhead[L]{\small\itshape Historical Jesus: An Evidentiary Reassessment}
\fancyhead[R]{\small\thepage}
\renewcommand{\headrulewidth}{0.4pt}

\usepackage[authoryear,sort&compress]{natbib}
\bibliographystyle{plainnat}
\usepackage{xcolor}
\usepackage[colorlinks=true,linkcolor=blue!60!black,
            citecolor=blue!60!black,urlcolor=blue!60!black]{hyperref}
\usepackage{booktabs}
\usepackage{enumitem}

\IfFileExists{csquotes.sty}%
  {\usepackage[autostyle,english=american]{csquotes}}%
  {\providecommand{\enquote}[1]{``##1''}}

\newcommand{\eg}{\textit{e.g.}}
\newcommand{\ie}{\textit{i.e.}}
\newcommand{\cf}{\textit{cf.}}
\newcommand{\circa}{\textit{c.}\,}
\newcommand{\grk}[1]{\textit{#1}}

\begin{document}

\begin{center}
  {\LARGE\bfseries The Historical Jesus:\\[4pt]
   An Evidentiary Reassessment}\\[18pt]
  {\large Stan Thomas}\\[4pt]
  {\normalsize Independent Researcher, St.\ George, Utah\\
   ORCID: 0000-0002-8828-7856}\\[12pt]
  {\normalsize Draft of \today}
\end{center}

\vspace{1em}

\begin{abstract}
\noindent
The historicity of Jesus of Nazareth is treated as settled in mainstream
New Testament scholarship.  This paper argues that the confidence is
unwarranted.  Drawing on a corpus-linguistic baseline study of 212
\textit{adelphos} tokens across six non-Christian Greek corpora, an
inscriptional survey of Greek religious-association usage of
\textit{adelphoi tou} [patron] (Harland 2005), a text-critical and
stylometric analysis of the Pauline passages historically adduced
against mythic-origin accounts, and a Bayesian sensitivity analysis
that implements source-dependency in the cumulative case, I argue
that (i)~the Pauline corpus describes a revelatory celestial figure
assembled from pre-Christian Jewish apocalyptic-mystical sources;
(ii)~the principal Pauline evidence traditionally cited against
mythic-origin scenarios, while non-trivial, is textually secure
and therefore not susceptible to interpolation arguments; (iii)~the
Gospels are not independent historical testimony but literary
constructions from Hebrew Bible templates; (iv)~the non-Christian
attestation is weaker than claimed and largely dependent on early
Christian informants; and (v)~under plausible assumptions the
posterior probability of historicity lands in open-question rather
than settled territory.  The responsible conclusion is that the
question is genuinely contested and that the scholarly consensus is
sustained in part by institutional and methodological inertia.
Full evidence and supplementary material are available in a
companion full-length treatment.

\vspace{0.5em}
\noindent\textbf{Keywords:} historical Jesus, mythicism, Pauline
epistles, corpus linguistics, \textit{adelphos tou kyriou}, Jewish
apocalyptic mysticism, Bayesian historiography.
\end{abstract}

%% ============================================================
\section{Introduction}
%% ============================================================

Does the evidence warrant the confidence of the historical-Jesus
consensus?  Mainstream New Testament scholarship treats the
historicity of Jesus as settled.  \citet{Ehrman2012} writes that
\enquote{virtually every competent scholar of antiquity, Christian
or non-Christian, agrees} that Jesus existed.  \citet{Dunn2003}
similarly: \enquote{today nearly all historians, whether Christians
or not, accept that Jesus existed.}  These claims are defensive
statements of consensus rather than evidentiary arguments, and
consensus about ancient persons is not evidence in the same sense
that primary attestation is.  The question this paper raises is
whether the primary evidence supports the confidence the consensus
projects.

The answer developed here is that it does not.  When each category
of evidence is examined on its own terms---without granting
historicity as a prior, without conflating consensus with
confirmation, and without suppressing the evidential weight of
absent evidence---the positive case is remarkably thin.  It reduces,
on close examination, to the interpretation of a single ambiguous
phrase in a single Pauline letter (Galatians~1:19), supported by
heuristic arguments about what is typical of religious movements.
The remaining evidence either fails textual scrutiny, derives from
the internal tradition of the movement itself, or is
indistinguishable from what we would expect if no specific
historical individual underlay the tradition.

The argument is not that Jesus definitely did not exist, and
intellectually honest mythic-origin scenarios do not claim as much.
Several passages in Paul---notably Galatians~4:4, Romans~1:3--4,
and the witness list in 1~Corinthians~15:3--8---pose genuine
challenges for mythic-origin accounts that must be acknowledged and
weighed.  The argument is that the evidence does not establish
historicity with the confidence the consensus claims, that a
well-constructed mythic-origin scenario drawing on Jewish
apocalyptic-mystical resources (not pagan mystery cults) is a
serious alternative that survives scrutiny of its weaknesses, and
that the consensus is sustained partly by the institutional
structure of the field rather than by evidential weight alone.

\paragraph{Scope of the present article.}
The full-length treatment of this argument (\textit{The Evidentiary
Status of the Historical Jesus: A Critical Assessment}, 84 pp.,
available at \url{https://github.com/sbthomas/historicity}) develops
the case at length, with two supplementary documents: a complete
colour-coded enumeration of every \textit{adelph-} family lexeme in
the seven undisputed Pauline letters, and a corpus-linguistic
baseline study of 212 random-sampled \textit{adelphos} tokens across
six non-Christian Greek corpora with an appendix of twenty-five
inscriptional cases of associative-community usage.  This article
presents the core argument in condensed form.  Evidence for which
space does not permit full development here is referenced to the
full-length treatment and its supplements.

%% ============================================================
\section{Methodological Preliminaries}
%% ============================================================

Four methodological commitments orient the argument.

\paragraph{Consensus is not evidence.}
Arguments from authority are heuristics, not independent evidential
contributions.  The epistemic value of a consensus depends on the
independence and quality of the underlying evidential judgements.
A consensus among scholars trained within the same paradigm, who
have absorbed the same default assumptions during professional
formation and who operate within institutions that presuppose the
conclusion in question, carries less evidential weight than a
consensus formed by independent evaluation of primary data.  This
paper treats the historicity consensus as a datum to be explained,
not as evidence to be weighed.

\paragraph{The Bayesian framework with source-dependence.}
Historical inference is inherently probabilistic.  \citet{Carrier2014}
applied a Bayesian framework systematically in \textit{On the
Historicity of Jesus}.  The framework itself is sound; the specific
point estimates Carrier assigns have been contested, as have the
independence assumptions underlying his multiplication of likelihood
ratios \citep{Vridar2014,Tucker2004}.  \citet{Tucker2004} makes the
critical observation that Bayesian historical inference is
particularly sensitive to the independence of contributing
testimonies: when textual-critical research shows that Gospel
attestations collapse to Markan or Q dependence, the effective
likelihood contribution is the contribution of Mark, not the
independent-multiplicative contribution of each synoptic Gospel
treated as independent.  A supplementary computational tool
(\textit{scripts/bayesian\_sensitivity.py} in the project repository)
implements this explicitly: under plausible point-estimate
likelihood ratios for each evidentiary strand and a neutral prior
$P(H)=0.5$, the dependency-aware posterior for historicity falls in
the range $0.23$--$0.39$---open-question territory by any reasonable
standard.  The exercise is a sensitivity analysis, not a definitive
computation.

\paragraph{Absence of evidence as evidence of absence.}
\enquote{Absence of evidence is not evidence of absence} is, in
general contexts, a truism that misrepresents the Bayesian point.
The correct principle is that absence of evidence \textit{is}
evidence of absence when the evidence would be expected to exist
under the hypothesis, with strength proportional to the
expectation \citep{Sober2009,Carrier2012}.  Where $n$ independent
channels each would produce evidence with probability $p$ under
hypothesis $H$ but probability $q < p$ under $\neg H$, the
cumulative likelihood ratio from silence is $(q/p)^n$, which
becomes small even when individual ratios are modest.  This
structure is developed below for the silence of Philo, Justus of
Tiberias, the Mishnah, and Roman administrative records.

\paragraph{The collapse of the criteria of authenticity.}
The traditional criteria of historical-Jesus research---embarrassment,
multiple attestation, dissimilarity, coherence---have been
subjected to devastating criticism from within mainstream
scholarship itself \citep{Allison2010,Keith2024}.  \citet{Keith2024}
and collaborators have explicitly argued that \enquote{the era of
the hard version of the criteria of authenticity is over}; the field
must move toward memory theory and reception history.  The paper
accepts this methodological consensus.  Its implication is that the
principal tools by which historicists have claimed to identify a
\enquote{historical core} beneath theological overlay are now
recognised as methodologically unreliable by the historicists
themselves.  If the criteria cannot reliably distinguish historical
material from theological invention, the Gospels cannot serve as
positive evidence for a specific historical figure.

\paragraph{Comparative perspective.}
Discussions of Jesus's historicity are routinely conducted as if the
question were unique.  It is not.  A class of comparable figures
whose putative historicity has been seriously questioned within
modern scholarship exists, and examining how those cases have been
resolved yields diagnostic criteria.  Pilate's existence was
settled by the 1961 Caesarea inscription \citep{Vardaman1962}.
Socrates is secured by multiple \textit{independent} contemporary
attestations---Plato, Xenophon, Aristophanes, Aristotle.
\textbf{Moses} is rejected by most modern Egyptologists on the
ground of archaeological silence in a place and period where
evidence would be expected \citep{FinkelsteinSilberman2001,ThompsonTL1999}.
\textbf{Lycurgus of Sparta} is treated as non-historical: late and
derivative literary sources, founder-myth conformity
\citep{Cartledge2001}.  \textbf{Romulus} is universally rejected as
historical despite a sophisticated literary tradition placing him in
a specific setting \citep{Beard2015,Wiseman1995}.
\textbf{Pythagoras}, \textbf{King Arthur}, \textbf{Homer}, and
\textbf{Laozi} are genuinely contested, with the consensus on each
oscillating between \enquote{probable historical kernel with
extensive legendary accretion} and \enquote{composite literary
construction} \citep{Riedweg2005,Higham2002,WestML2011,GrahamA1986}.

Applying the diagnostic criteria to Jesus: he lacks the direct
contemporary physical attestation of Pilate, he lacks the multiple
independent contemporary literary attestation of Socrates, and his
Gospel biographical sources are literary-critically constructed
from prior templates (as with Lycurgus, Romulus, and Laozi).  His
profile most closely matches the \textbf{Pythagoras/Arthur/Homer/Laozi}
category of genuinely contested figures, not the Pilate/Socrates
category of secured figures.  The methodological implication is
that the field's confidence in historicity is incommensurate with
the evidence by the standards of comparable existence-questions.

%% ============================================================
\section{The Pauline Evidence}
%% ============================================================

The letters attributed to Paul constitute the earliest stratum of
surviving Christian writing.  If a historical Jesus existed and was
known to Paul's contemporaries, this is where we would expect to
find the clearest traces.  The Pauline corpus is therefore the most
important body of evidence in the historicity debate.

%% ------------------------------------------------------------
\subsection{The Character of Paul's Jesus}
%% ------------------------------------------------------------

The central observation, developed in detail by
\citet{Doherty1999,Doherty2009}: Paul's Jesus lacks the attributes
of a recently-living historical person.  Paul never names Jesus's
parents, birthplace, or hometown.  He provides no physical
description.  He narrates no events from Jesus's earthly
life---no ministry, no healings, no exorcisms, no parables, no
conflicts with authorities, no trial, no specific date or location
of execution.  He does not identify Judas as a betrayer or name
Pilate.  He does not place the crucifixion at Golgotha, in
Jerusalem, or indeed anywhere on earth.

What Paul does describe is a cosmic soteriological drama.  Jesus is
a pre-existent divine figure who \enquote{emptied himself} and took
on the \enquote{form of a servant} (Philippians~2:5--11---itself
likely a pre-Pauline hymn), who was \enquote{declared to be Son of
God in power according to the Spirit of holiness by his
resurrection from the dead} (Romans~1:4), who was \enquote{handed
over} (\grk{paradidomi}, a term routinely translated
\enquote{betrayed} in 1~Corinthians~11:23 to harmonise with the
Gospel narrative but whose primary meaning is
\enquote{delivered up}), and who died \enquote{in accordance with
the scriptures} and was raised \enquote{on the third day in
accordance with the scriptures} (1~Corinthians~15:3--4).

Paul explicitly tells us how he came to know this Christ:
\enquote{I did not receive [the gospel] from any human source nor
was I taught it, but I received it through a revelation
(\grk{apokalypsis}) of Jesus Christ} (Galatians~1:12);
\enquote{when God\ldots was pleased to reveal his Son in me}
(Galatians~1:16); \enquote{I went up in response to a revelation}
(Galatians~2:2); \enquote{I know a person in Christ who fourteen
years ago was caught up to the third heaven\ldots caught up into
Paradise and heard things that are not to be told}
(2~Corinthians~12:2--4).  This is the epistemic profile of a Jewish
mystical-apocalyptic visionary, not of a disciple of a recently-
living teacher.  We return to this point below.

%% ------------------------------------------------------------
\subsection{\textit{Adelphos tou Kyriou}: \enquote{Brother of the Lord}}
%% ------------------------------------------------------------

If the foregoing is correct, the principal historicist evidence
from the Pauline corpus is a single passage: Galatians~1:19.
Paul, describing his first visit to Jerusalem, writes that he saw
no other apostle \enquote{except James, \grk{ton adelphon tou
kyriou}---the brother of the Lord.}  The historicist reading is
straightforward: James was the biological brother of Jesus.
\citet{Ehrman2012} considers this passage virtually decisive.

The passage deserves careful scrutiny because it is being asked to
carry the weight of the historicist case.  Three considerations
complicate the biological-sibling reading.

\paragraph{Paul's own distributional usage.}
The seven undisputed Pauline letters contain exactly 121 occurrences
of the \grk{adelph-} family of lexemes (\grk{adelphos},
\grk{adelph\=e}, \grk{pseudadelphos}, \grk{philadelphia}).  Of
these, 117 are unambiguously community/fellow-believer uses, 3 are
disputed (1~Cor~9:5 \grk{hoi adelphoi tou kyriou}, Gal~1:19 itself,
Rom~16:15 Nereus and his sister), and 1 is ethnic/covenantal
(Rom~9:3 \grk{adelph\^on mou kata sarka}, fellow-Israelites).
\textbf{Zero} Pauline tokens unambiguously denote a biological
sibling outside the two disputed \grk{tou kyriou} passages.  This
is not a selective sample; it is the complete corpus, colour-coded
in the project's Pauline supplement.

The historicist reply is that \grk{adelphos tou kyriou} is a
\textit{specifically modified} construction and that the 117
generic uses of plain \grk{adelphos} do not govern the kinship
reading of the genitive phrase.  That reply assumes the existence
of a separate biological-sibling usage from which the modified
form is drawn.  In Paul, there is no such usage.  The biological
reading of Galatians~1:19 requires Paul to have used
\grk{adelphos} in a unique sense in this one phrase---a lexical
register Paul exhibits nowhere else in his authentic corpus.

\paragraph{The parallel in 1~Corinthians~9:5.}
The decisive test case is 1~Corinthians~9:5: \enquote{Do we not
have the right to be accompanied by a believing wife, as do the
other apostles and \grk{hoi adelphoi tou kyriou} and Cephas?}
Paul is listing categories of authoritative figures:
(i)~the other apostles, (ii)~the brothers of the Lord,
(iii)~Cephas.  Three distinct groups or individuals, cited in
parallel as recognised authority classes.  \enquote{The brothers
of the Lord} is \textit{plural} and functions as a class label.

If this phrase means \enquote{biological siblings of Jesus}, then
there exists a group of biological siblings sufficiently prominent
to be cited as a recognised authority class alongside the
apostolic college and Peter individually.  The simpler reading is
available: \enquote{brothers of the Lord} names a designated group
or rank within the movement---an inner circle, a faction, or a
class---distinguished from \enquote{apostles} and from individual
leaders like Cephas.

\paragraph{The corpus-linguistic baseline.}
A corpus-linguistic study (full methodology in the supplementary
\textit{baseline} document) random-samples 212 \grk{adelphos}
tokens from six non-Christian contemporary-or-earlier Greek
corpora: papyri (100~BCE--70~CE), the Septuagint, Philo of
Alexandria, Josephus, pre-Pauline literary (Polybius, Diodorus,
Strabo), and Hellenistic literary (Plutarch, Epictetus, Dio
Chrysostom).  Classification into A~(community / voluntary
association), B~(biological sibling), C~(disputed), D~(ethnic /
covenantal), E~(literary metaphor) yields:

\begin{center}
\small
\begin{tabular}{lrrrrr}
\toprule
Corpus & A\% & B\% & C\% & D\% & E\% \\
\midrule
Papyri (pre-70 CE), $n{=}35$ & 0.0 & 100.0 & 0.0 & 0.0 & 0.0 \\
LXX, $n{=}35$                & 0.0 & 68.6  & 5.7 & 25.7 & 0.0 \\
Josephus, $n{=}35$           & 0.0 & 60.0  & 5.7 & 34.3 & 0.0 \\
Philo, $n{=}35$              & 2.9 & 57.1  & 2.9 & 5.7  & 31.4 \\
Hellenistic lit., $n{=}36$    & 0.0 & 80.6  & 0.0 & 2.8  & 16.7 \\
Pre-Pauline lit., $n{=}36$   & 2.8 & 94.4  & 0.0 & 0.0  & 2.8 \\
\midrule
\textbf{Aggregate}, $N{=}212$ & \textbf{0.9} & \textbf{76.9} & 2.4 & 11.3 & 8.5 \\
\midrule
\textbf{Paul}, $N{=}121$     & \textbf{96.7} & \textbf{0.0} & 2.5 & 0.8 & 0.0 \\
\bottomrule
\end{tabular}
\end{center}

The nominal ratio of Pauline-to-baseline A-rate is approximately
103~to~1; Wilson 95\% confidence intervals do not overlap.  This
figure is a \textit{lower bound} of the ratio: the sampled baseline
excludes the Greek inscriptional register, where, as
\citet{Harland2005} and the Kloppenborg-Ascough \citep{KloppenborgAscough2011}
corpus document, voluntary religious-association usage of
\grk{adelphoi} concentrates.  Factoring inscriptions in would raise
the baseline A-rate to roughly 5--15\% and compress the ratio to
roughly 10--20$\times$.  The qualitative finding---that Paul's
usage is distinctive against ambient Greek---survives this
correction.  The quantitative claim of 103$\times$ does not, and
should not be cited in isolation.

What the corrected analysis reveals is interesting in its own
right.  Paul's usage is \textit{not} distinctive against the
Greek religious-associative sub-register, where community-title
\grk{adelphoi} with genitive-of-patron is well attested.  Twenty-
five inscriptional cases documented in the supplementary
\textit{baseline} document (principally the \enquote{adopted
brothers of Theos Hypsistos} at Tanais, c.~155--240~CE) establish
the specific construction \grk{adelphoi tou} [patron figure] as
formally normal in Greek voluntary religious associations of the
period.  The Pauline \grk{adelphoi tou kyriou} instantiates that
convention; it does not require it, but it does not depart from
it either.

\paragraph{Assessment.}
When the three considerations combine---no biological-sibling
register in Paul's own corpus, the class-label parallel in
1~Cor~9:5, and the well-attested Greek associative-community
convention---the community-title reading of Galatians~1:19 is the
more parsimonious reading.  The historicist reading is possible
but requires Paul to depart from his own attested usage, from the
parallel within his own corpus, and from a convention attested in
the religious-associative sub-register within which he is writing.

Galatians~1:19 is \textit{compatible with} historicity; it does
not \textit{require} it.

%% ------------------------------------------------------------
\subsection{The Passages That Do Challenge Mythic-Origin Scenarios}
%% ------------------------------------------------------------

Several Pauline passages pose genuine difficulties for mythic-
origin scenarios and must be acknowledged rather than dismissed.

\paragraph{Romans~1:3.}
\enquote{descended from the seed of David according to the flesh}
(\grk{tou genomenou ek spermatos Dauid kata sarka}).  Davidic
descent is a genealogical claim.  \citet{Novenson2012,Novenson2017}
has documented that messianic language in Second Temple Judaism
operates within constraints that make \enquote{messiah} a human-
historical designation, not a detachable cosmic abstraction.  The
mythic-origin response runs along three lines: first, that
Romans~1:3--4 is widely recognised as a pre-Pauline creedal
fragment that Paul quotes rather than composes \citep{Jewett2007};
second, that \grk{kata sarka} (\enquote{according to the flesh})
is a technical Pauline phrase that can locate an event in the
sublunar realm without committing to a specific earthly biography
\citep{Carrier2014}; third, that Davidic-descent claims could arise
through scriptural exegesis (reading Jer~23:5, Ps~89, 2~Sam~7 as
requiring a Davidic lineage for any messianic figure) rather than
through biographical memory.  None of these responses fully
dissolves the tension.  The honest concession is that Rom~1:3--4
is weak evidence for historicity that the mythic-origin scenario
must accommodate.

\paragraph{Galatians~4:4.}
\enquote{born of a woman, born under the law}.  A bare
incarnation-into-Jewishness claim that could reflect memory of a
historical Jewish Jesus or could reflect the heavenly mediator's
descent through an earthly birth into the covenantal community.
Like Rom~1:3, compatible with both hypotheses but more naturally
historicist.

\paragraph{1~Corinthians~15:3--8.}
The resurrection creed with witness list.  Verses~3--5 are a
universally-acknowledged pre-Pauline formula; verses~6--8 extend to
a specific list of resurrection appearances (Cephas, the Twelve,
five hundred brothers, James, all the apostles, Paul himself).
The specificity of the witness list, especially the \enquote{over
five hundred brothers} figure, is the hardest Pauline datum for
mythic-origin scenarios.  The mythic-origin response locates the
witness list within the established Second Temple pattern of
\textit{post-mortem visionary experience attributed to multiple
members of a sectarian community}, the paradigm example of which
is the Qumran Self-Glorification Hymn (4Q491c) where an unnamed
speaker claims heavenly ascent and reception by the angelic council
\citep{Collins1995,Knohl2000}.  Visionary experience of a heavenly
figure by multiple community members is attested in pre-Christian
Jewish mystical-apocalyptic contexts.  This does not fully
dissolve the tension either.

\paragraph{Textual and stylometric stability.}
A text-critical and preliminary stylometric analysis (full details
in the full-length paper, §3.1, and in
\textit{scripts/stylometry.py}) establishes that these contested
passages are \textit{textually secure}.  They are attested in
\textbf{P46} (the Chester Beatty Pauline papyrus, \circa 200~CE),
in all major subsequent codices, and are present in Marcion's
reconstructed text where reconstruction is possible.  Stylometric
distance-from-baseline, which correctly identifies known
interpolations (1~Thess~2:14--16 and 1~Cor~14:34--35) as outliers,
does \textit{not} flag the historicity-relevant passages as
interpolation candidates: Galatians~1:18--20 in context has a
stylometric distance from the Pauline baseline of 0.09, less
distinctive than established interpolations ($\approx 0.10$) and
less distinctive than the recognised pre-Pauline creedal formulae
($\approx 0.17$).  The mythic-origin response to these passages
\textit{cannot} therefore rely on text-critical or stylometric
excision.  It must rely on semantic and contextual analysis, which
is what the paper provides.

\paragraph{Net Pauline picture.}
The Pauline evidence is a mixed landscape.  The \grk{adelphos tou
kyriou} passages, when analysed in their full Pauline distribution
and against Greek associative-community convention, do not secure
the historicist reading.  Paul's broader presentation of Jesus as
a revelatory celestial figure known through \grk{apokalypsis} and
scripture does not require a historical subject to explain the
data.  The three conceded-difficulty passages
(Rom~1:3--4, Gal~4:4, 1~Cor~15:3--8) admit mythic-origin readings
that require interpretive work the historicist readings do not;
this asymmetry is real and must be acknowledged.  The full
analysis thus lands between the confident-historicist and
confident-mythicist positions: the question is genuinely open.

%% ============================================================
\section{The Jewish-Apocalyptic Matrix}
%% ============================================================

An older popularisation of the mythic-origin hypothesis, associated
with \citet{Frazer1890}, located its sources primarily in pagan
Mediterranean mystery cults---Osiris, Attis, Adonis, Dionysus---and
is what many readers think of when they encounter the phrase
\enquote{mythicist hypothesis}.  That framing is the weakest
version of the argument.  As \citet{Fredriksen2017,Fredriksen2018}
insists, Paul is a Second Temple Jewish apocalyptic thinker
operating within Jewish categories, not a syncretist blending in
pagan myth.  The honest mythic-origin argument is that the Jesus
tradition is constructed from the scriptural, apocalyptic, and
mystical resources of Second Temple Judaism itself, with Hellenistic
literary conventions providing secondary ambient shaping.

\paragraph{Pre-existent heavenly mediators.}
\citet{Boyarin2012} traces an attested \enquote{binitarian} strand
through the Danielic Son of Man (Daniel~7:13--14), the
\textit{Similitudes of Enoch} (1~Enoch~37--71), the Melchizedek
traditions at Qumran (11QMelch), angelomorphic mediator figures
across Second Temple apocalyptic literature, and the Logos of
Philo of Alexandria.  Paul's Christ is not an exotic Hellenistic
import into Judaism; he is a recognisable member of the Jewish
pre-existent-mediator family.  The Philippians~2:5--11 descent-
and-exaltation arc maps onto Second Temple patterns of
angelomorphic humiliation and vindication (the Self-Glorification
Hymn at Qumran, 4Q491c), not onto the death-and-rebirth cycle of
Attis.

\paragraph{Philo's Logos.}
\citet{Sterling2001} documents Philo's Logos theology in detail.
The Logos is pre-existent, functions as God's agent in creation,
is called \enquote{son of God} and \enquote{firstborn}, is the
mediator through whom humans approach God, and is sometimes
described as \enquote{a second god} (\textit{deuteros theos},
\textit{Questions on Genesis} 2.62).  Philo wrote in Alexandria in
the same decades as Paul's formation.  The conceptual inventory
from which early Christological language could be assembled was
already available in Jewish-Hellenistic form; no pagan borrowing
is required.

\paragraph{The revelatory mechanism.}
Second Temple Judaism possessed a specific, well-attested set of
\textit{procedures} for generating religious knowledge about
heavenly realities: visionary ascent through the heavens plus
scriptural exegesis decoding the Hebrew Bible as encoding
cosmic mysteries.  \citet{Scholem1941}, \citet{Gruenwald1980},
\citet{Rowland1982}, and \citet{SchaferP2009} have documented this
tradition across 1~Enoch, the Apocalypse of Abraham, the Ascension
of Isaiah, the Qumran Self-Glorification Hymn, and the Merkabah
and Hekhalot literatures.  Paul's explicit claims about his
Christ-knowledge---received by \grk{apokalypsis} (Gal~1:12),
revealed \enquote{in me} (Gal~1:16), received through
\enquote{caught up to the third heaven}
(2~Cor~12:2--4)---place him squarely within this Jewish mystical-
apocalyptic procedural framework.  \citet{Segal1990} argued in
detail that Paul's self-description matches the profile of a
Jewish mystical visionary.

None of the texts in this tradition requires a recent historical
instantiation of the heavenly figure to explain the revelatory
knowledge.  1~Enoch's Son of Man is a named, hidden, pre-existent
heavenly figure who will appear at the end of days---not the memory
of a recently-living person.  The Self-Glorification Hymn's
enthroned speaker is a visionary claim, not a biographical report.
The Christian community could have come to know its Christ by the
same procedural route---scriptural exegesis plus visionary
revelation---that comparable Second Temple communities used to
come to know theirs.

\paragraph{Crucifixion and Deuteronomy~21:23.}
The specific mode of death is not well explained by pagan dying-
and-rising cults (where divine deaths take varied forms).  It is
well explained by scriptural logic.  Deuteronomy~21:23:
\enquote{anyone hung on a tree is under God's curse.}  Paul
explicitly cites this passage in Galatians~3:13: \enquote{Christ
redeemed us from the curse of the law by becoming a curse for
us---for it is written, `Cursed is everyone who hangs on a
tree.'\,}  This is the engine of Paul's substitutionary-atonement
soteriology.  For the theology to work, the mode of death must be
\textit{scripturally cursed}.  Stoning, burning, beheading---none
of these carry the specific scriptural curse that hanging on a
tree does.  Crucifixion was chosen not \textit{despite} its
scandalous character but \textit{because} it uniquely satisfied
the scriptural requirement for a substitutionary-curse theology.

\paragraph{Significance.}
Every major element of Paul's Christology---pre-existence, descent,
vicarious suffering, resurrection, exaltation, cosmic lordship---
is derivable from Jewish scriptural and apocalyptic sources
available in the first century.  This is \textit{a fortiori} the
safer argumentative ground: the Jewish sources are textually
secure (the Qumran material was effectively sealed by 70~CE and
shows no Christian influence), whereas the pagan parallels face
genuine scholarly contestation about dating, comparability, and
direction of influence.  The honest mythic-origin scenario does
not require recourse to Attis or Osiris.  Second Temple Judaism
supplies the full theological machinery.

%% ============================================================
\section{The Gospel Evidence}
%% ============================================================

The canonical Gospels are routinely treated as evidence for
historicity, both by the consensus and in popular understanding.
The question is whether that treatment is warranted.

It is uncontroversial in mainstream scholarship that the Gospels
are not eyewitness accounts \citep{Ehrman2009}.  They are
anonymous works (the names Mark, Matthew, Luke, John are second-
century attributions), written in Greek by educated authors,
forty to seventy years after the events they purport to describe,
in communities far removed from rural Galilee.  They are
theological documents built from identifiable literary templates.

\paragraph{Scriptural construction of the passion narrative.}
The passion narrative is constructed, point by point, from Hebrew
Bible texts.  \textbf{Psalm~22:} \enquote{My God, my God, why have
you forsaken me?} (Mk~15:34 / Ps~22:1); \enquote{They divide my
garments} (Mk~15:24 / Ps~22:18); \enquote{All who see me mock me}
(Mk~15:29 / Ps~22:7); \enquote{They pierce my hands and my feet}
(Ps~22:16, reflected in the crucifixion itself).
\textbf{Isaiah~53:} \enquote{He was oppressed and afflicted, yet
he did not open his mouth} (Is~53:7; Mk~15:5); \enquote{He was
numbered with the transgressors} (Is~53:12; Mk~15:27);
\enquote{He bore the sin of many} (Is~53:12).
\textbf{Wisdom of Solomon 2--5} provides the structural parallel
of persecuted-righteous-one persecuted and vindicated
\citep{Nickelsburg1972}.  \textbf{Zechariah:} the entry on a
donkey (Zech~9:9 / Mk~11:1--7); thirty pieces of silver
(Zech~11:12--13 / Mt~26:15); the shepherd struck, the sheep
scattered (Zech~13:7 / Mk~14:27).  \textbf{Daniel~7:} the Son of
Man coming on clouds (Dn~7:13--14 / Mk~14:62).

The same pattern extends across the Gospel narrative arc: birth
narratives built from Moses, Samuel, Isaac typology; ministry
episodes recapitulating Elijah-Elisha; the ascension drawing on
2~Kings~2:11 and Dn~7:13--14.  When every major episode can be
derived from a specific scriptural source, the hypothesis that
historical events \textit{preceded} the scriptural interpretation
becomes less parsimonious than the hypothesis that scriptural
sources \textit{generated} the narrative.  The tradition reads
as a systematic programme of scriptural fulfilment---exactly how
the Gospel authors themselves present it (Matthew's formula
citations).

\paragraph{What this establishes and what it does not.}
It is important to be precise about the logical step the argument
makes, because a historicist critic such as \citet{Allison2010}
would correctly insist that scriptural templating of a narrative
does not entail the non-existence of the subject.  Historical
figures have been described in antiquity through scriptural and
typological grammars.  Jewish historiography routinely employs
Hebrew Bible typology to narrate the lives of real people
(Maccabean heroes as new Judges, exilic returnees as new Exodus).
That the Gospel passion is constructed from Ps~22 and Is~53 does
not on its own prove that no crucifixion of Jesus occurred.

The conclusion the argument supports is narrower but sufficient.
It is that the Gospel passion narrative is not \textit{independent
historical testimony}.  Its content is demonstrably derived from
Hebrew Bible templates rather than from witness testimony, oral
tradition traced to an event, or documentary record.  The Gospels
cannot therefore serve as evidence \textit{for} the historicity
of the crucifixion; they can serve only as evidence for the fact
that late first-century communities constructed theologically
meaningful narratives of a crucifixion using available scriptural
materials.

Whether any historical event lies beneath those constructions must
be established, if at all, from sources that are not themselves
scriptural derivations.  Paul might be such a source, and §3 has
argued that the Pauline description of the crucifixion is
generic, soteriological, and non-biographical.  Josephus might be
such a source, and §5 will argue his attestation is weaker than
usually claimed.  The argument is cumulative: templating
neutralises the Gospels as independent testimony; the remaining
strata must carry the full historicist burden.  They cannot carry
that burden without assistance from the Gospels they are being
asked to corroborate.

%% ============================================================
\section{Non-Christian Attestation}
%% ============================================================

\paragraph{Josephus \textit{Antiquities} 18.63--64 (the
\textit{Testimonium Flavianum}).}
The passage as transmitted contains claims no non-Christian Jewish
historian of the first century would plausibly make
(\enquote{He was the Christ}; \enquote{when Pilate had condemned
him to the cross\ldots he appeared to them alive again on the
third day}).  The consensus position, following
\citet{Meier1991,Ehrman2012}, is that the \textit{Testimonium} is
partially interpolated; \citet{Olson1999} has argued that the full
text is an Eusebian composition (\textit{c.}~324~CE).
\citet{Carrier2014} adopts the full-interpolation position.
Either way, the \textit{Testimonium}'s evidential weight for
historicity is weak.  The authentic core, if any, is at most the
observation that a Christian community traced to a figure called
Jesus existed by the late first century---which neither side of
the debate contests.

\paragraph{Josephus \textit{Antiquities} 20.200 (the James
passage).}
\enquote{James, the brother of Jesus who was called Christ.}
This passage is accepted by most textual scholars as substantially
authentic and is regarded as the strongest non-Pauline attestation
for historicity.  Its evidential weight is nonetheless limited.
Josephus is writing \circa 93--94~CE, sixty years after any
supposed crucifixion, at a time when Christian communities were
already spread across the eastern Mediterranean and their
narrative about their founder was already in circulation.  His
passage attests what Christian communities by that date claimed
about their founder, not independently that Jesus was a
first-century historical individual.  The structure is identical
to Tacitus's reference (below): a mid-to-late first-century source
reports the Christian community's own claim.  Under the
mythic-origin scenario this is exactly what we would expect;
under the historicist scenario it is exactly what we would expect;
the evidence does not discriminate.

\paragraph{Tacitus \textit{Annals}~15.44.}
Written \circa 116~CE, describing Nero's persecution of Christians
(64~CE) and referencing \enquote{Christus, from whom the name had
its origin, suffered the extreme penalty during the reign of
Tiberius at the hands of one of our procurators, Pontius
Pilatus.}  The passage is almost certainly dependent on Christian
informants or on Christian tradition already circulating in Rome;
Tacitus is not citing Roman administrative records (which would
not have preserved the execution of a Galilean preacher a century
earlier) but the Christian community's self-description as
Tacitus received it.  Like Josephus 20.200, the passage confirms
early circulation of a historicised narrative without independently
supporting the narrative itself.

\paragraph{Rabbinic and Mandaean traditions.}
The Talmudic Yeshu passages (b.Sanhedrin~43a, 107b; b.Sotah~47a;
b.Shabbat~104b), systematically surveyed by \citet{SchaferP2007},
and the \textit{Toledoth Yeshu} compilations
\citep{MeersonSchafer2014}, represent polemical Jewish response to
Christian claims rather than independent attestation.  Similarly,
the Mandaean tradition---a surviving Mesopotamian community
centred on John the Baptist with its own Jesus-references in the
\textit{Ginza Rabba} and \textit{Book of John}
\citep{Rudolph1978,BuckleyMandaeans2002,Lupieri2002}---treats Jesus
as a polemically-reshaped figure whose existence is presupposed
but whose significance is inverted.  These are interesting data
for the diversity of the Second Temple religious environment but
are too late and too dependent on the Christian narrative to count
as independent attestation of historicity.

\paragraph{Pliny, Suetonius, and minor references.}
\textit{Pliny the Younger}'s letter to Trajan (\circa 112~CE)
confirms the existence of Christians in Bithynia; it tells us
nothing about a first-century historical figure.  \textit{Suetonius}'s
reference to \enquote{Chrestus} (\textit{Life of Claudius} 25.4)
may not be to Christ at all; even if it is, it attests Jewish
unrest in Rome \circa 49~CE, not the historicity of Jesus.

\paragraph{Net non-Christian picture.}
The non-Christian attestation is weaker than typically claimed.
What it establishes uncontestedly is that by the late first and
early second centuries the historicised Christian narrative was
in wide circulation across the eastern Mediterranean and had
reached Roman awareness.  Both historicist and mythic-origin
scenarios must accommodate this datum; neither is favoured by it.

%% ============================================================
\section{The Argument from Silence}
%% ============================================================

Section~2 established the formal principle: absence of expected
evidence is evidence of absence, with strength proportional to the
expectation.  Four documentary channels are directly relevant.

\paragraph{Philo of Alexandria} (\circa 20~BCE--\circa 50~CE): an
exact contemporary of the supposed Jesus, a Hellenistic Jewish
philosopher with an enormous surviving corpus, who wrote
extensively about Pilate, about Jewish messianic expectation, about
Wisdom and Logos theology---the precise intellectual milieu from
which Christianity emerged---and who mentions Jesus nowhere.  The
silence is moderated by Philo's Alexandrian location, but it is
not eliminated.

\paragraph{Justus of Tiberias} (late first century): a historian of
Galilee---Jesus's supposed home region---whose silence on Jesus
was explicitly remarked by Photius in the ninth century when he
read the now-lost \textit{History}.  A native Galilean historian's
silence about a Galilean religious phenomenon substantial enough
to generate a Mediterranean-wide movement within one generation is
among the strongest individual silences in the case.

\paragraph{The Mishnaic tradition.}
The Mishnah, codified \circa 200~CE from earlier oral traditions,
is a legal code rather than a chronicle.  Its silence on Jesus is
a weaker datum because the Mishnaic genre does not demand reference
to individual historical figures.

\paragraph{Roman administrative records.}
No direct record of Jesus's trial or execution survives.  The
silence is moderated by the general loss of provincial records
from the period and by the low-profile character any peasant
preacher's execution would have had within Roman administration.

\paragraph{Cumulative assessment.}
The cumulative likelihood ratio from silence is difficult to
quantify precisely.  The qualitative conclusion is defensible:
the pattern of silence across multiple independent traditions is
\textit{more consistent with} mythic-origin scenarios than with
historicist scenarios, because mythic-origin scenarios predict
silence while historicist scenarios must explain it through
auxiliary hypotheses about obscurity.  The cumulative weight is
real but modest, not decisive; it is one strand of a cumulative
case, not a standalone proof.

%% ============================================================
\section{Conclusion}
%% ============================================================

The evidence for the historicity of Jesus of Nazareth has been
examined across every available category.  The findings may be
summarised as follows.

The Pauline corpus describes a revelatory celestial figure known
through \grk{apokalypsis} and scriptural exegesis.  Paul's
discourse is intelligible within the Second Temple Jewish mystical-
apocalyptic tradition without requiring a recent historical
subject.  The \grk{adelphos tou kyriou} passages
(Gal~1:19, 1~Cor~9:5)---the principal historicist evidence---are
more parsimoniously read as community-title than as biological
kinship, given Paul's own corpus-wide distribution (117~community
uses, 0~unambiguous biological uses outside the two contested
cases) and the well-attested Greek associative-community use of
\grk{adelphoi tou} [patron figure] documented by
\citet{Harland2005}.  The contested passages that do pose
difficulties for mythic-origin scenarios
(Rom~1:3--4, Gal~4:4, 1~Cor~15:3--8) are textually secure and
admit mythic-origin readings, though with interpretive work the
historicist readings do not require.  This asymmetry is honestly
acknowledged.

The Gospel narratives are demonstrably constructed from Hebrew
Bible templates.  This establishes that the Gospels are \textit{not
independent historical testimony}, not that they disprove the
existence of a subject.  The evidentiary burden therefore shifts
to non-Gospel sources, which are individually weaker than the
consensus projects.

The non-Christian attestation (Josephus, Tacitus, rabbinic,
Mandaean, Pliny, Suetonius) confirms the early circulation of a
Christian narrative about a historical founder but does not
independently confirm the narrative.  Both hypotheses accommodate
this datum.  The cumulative silence of contemporaneous sources
(Philo, Justus of Tiberias, the Mishnah, Roman administrative
records) constitutes qualitatively real disconfirming evidence
whose aggregate weight is modest but real.

Every narrative element attributed to the earliest Christian
community's Christ---pre-existence, descent, vicarious suffering,
resurrection, exaltation, cosmic lordship, crucifixion under
scriptural curse---is derivable from pre-Christian Jewish
apocalyptic-mystical resources (the Danielic Son of Man, Philo's
Logos, angelomorphic mediators, the Suffering Servant,
Dt~21:23 in combination with scriptural exegesis).  The pagan-
mystery-cult framing of the older mythicist literature is not the
argument this paper defends.

Placed in the landscape of comparable existence-questions,
Jesus's evidentiary profile most closely matches the contested-
figures category (Pythagoras, Homer, Arthur, Laozi)---probable
historical kernel with massive legendary elaboration, or community
construction to satisfy founder-myth requirements, with the
evidence insufficient to discriminate---rather than the secured-
figures category (Pilate, Socrates) or the rejected category
(Moses, Lycurgus, Romulus).  Under plausible Bayesian assumptions,
the posterior probability of historicity lands in the range
0.23--0.39 given a neutral prior, and materially below the near-
certainty the consensus projects even with a strong historicist
prior.

The responsible conclusion is that the evidence does not establish
the historicity of Jesus with the confidence the consensus claims,
that the mythic-origin hypothesis in its Jewish-apocalyptic form
is a serious alternative that survives scrutiny of its weaknesses,
and that the scholarly consensus is sustained partly by the
institutional and methodological structure of the field.  This
paper does not argue that Jesus definitely did not exist; it does
not argue that mythic-origin has been proven more probable than
historicity; it argues that the evidence is insufficient to
support near-certainty either way.  The question should be treated
as genuinely open, and if the field cannot bring itself to treat
it as open---if the mere posing of the question continues to be
treated as a mark of incompetence rather than as a legitimate
scholarly inquiry---the field has a problem that no amount of
consensus can resolve.

\vspace{1em}

\noindent\rule{0.4\textwidth}{0.4pt}\par
\vspace{0.5em}
\small\noindent\textit{Supplementary material.}  Complete data and
methodology in the full-length companion paper \textit{The
Evidentiary Status of the Historical Jesus: A Critical Assessment}
(84~pp.) and two supplements: a colour-coded enumeration of every
\grk{adelph-} lexeme in the seven undisputed Pauline letters, and
a corpus-linguistic baseline study of 212 random-sampled
\grk{adelphos} tokens across six non-Christian Greek corpora with
inscriptional appendix.  All artifacts at
\url{https://github.com/sbthomas/historicity}.

\noindent\textit{Acknowledgments.}  This paper benefited from
extensive critical dialogue conducted with Claude (Anthropic),
which served as adversarial interlocutor and peer reviewer
throughout.  The arguments and conclusions are the author's own.

\bibliography{historicity}

\end{document}
